\scanpage{3}
\uline{Rappels de théorie des groupes} \uline{et sur espaces homogènes}

Groupe (loi comp. assoc. unitaire avec inverses)

hom de groupes (on a néc $f(e)=e$, $f(x^{-1})=f(x)^{-1}$)

opération d'un groupe $G$ sur un ens (ens. à groupe d'opérateurs)

opération transitive ($\forall x\mathord{.}y\in E,\ \exists g\in G,\ y=gx$) -- espace homogène

On exige \deleted{que}\inserted{de plus} $E\ne\varnothing$)
\[
\bigl(\text{couples }(E,x),\ E\ \text{esp. hom.},\ x\in E\bigr)
\Longleftrightarrow
\begin{aligned}&\text{donnée d'un ss-grpe}\\&H\ \text{de }G\end{aligned}
\]

c'est en fait une ``équivalence de catégories''

% slanted
\margin{classes d'iso. d'esp. hom. ss $G$ $\simeq$ classes de conjugaison de ss-grpes}
{\small
\diagram{\begin{tikzcd}[ampersand replacement=\&,row sep=0.6em,column sep=0.8em]
(E,x) \arrow[dd,"f"] \& \& \& {}\\
{} \& (f(x)=x') \& \Longleftrightarrow \& H\subset H'\\
(E',x') \& {\begin{aligned}&\text{\uline{NB} }f\ \text{(comm. à }G\text{) connu,}\\&\text{quand on \unsure{connait|connaît}}\\& f(x),\ \text{\uline{si} }G\ \text{trans. sur }E\end{aligned}} \& \& {}
\end{tikzcd}}}
{\small
\[
E\sim E'\quad\text{sss}\quad H,H'\ \text{conjugués}\quad\big/\quad
\exists E\to E'\quad\text{sss}\quad
\begin{aligned}&H\ \text{conjugué à un}\\[-2pt]&\text{ss-grpe de }H'\end{aligned}
\]}

\uline{Torseur sous $G$} :
{\small
\diagram{\begin{tikzcd}[ampersand replacement=\&,row sep=1.25em,column sep=1.5em]
{\left\{\begin{aligned}&\forall x,y\in E,\quad\exists 1\ g\in G,\quad gx=y\\& E\ne\varnothing\end{aligned}\right\}}
\arrow[r,Leftrightarrow] \arrow[d,Leftrightarrow] \&
{\left\{\begin{aligned}&E\ \text{\unsure{espace|esp.} hom.,}\\&\text{stabilisateur}\\&\text{\unsure{de|des} pts réd.}\\&\text{à }\{e\}\end{aligned}\right\}}\\
{\left\{E\simeq G_{s}\right\}} \&
\end{tikzcd}}}

\deleted{\unsure{Les|Le}}

Les torseurs sous $G$ sont isomorphes (non canoniquement)

Les torseurs sous $G$ sont les ``plus gros'' espaces homogènes ss $G$.

le groupe des autom. d'un torseur $P$ sous $G$ est $\sim G$ (non canonique) : le choix de $x\in P$ définit un iso. $G\overset{u_x}{\simeq}\operatorname{Aut}_{G}P$, si on remplace $x$ par \inserted{$y=$}$gx$, cet iso. est remplacé par
\[
\boxed{u_y=u_{gx}=\deleted{\unsure{\operatorname{int}(g)\circ|\operatorname{int}g\circ}}u_x\circ\,?\,\operatorname{int}(s^{-1})}
\]
{\small
\[
\begin{aligned}
\bigl[\,u_x(g).\,x=gx\qquad &u_y(g)y=gy=gsx\unsure{\deleted{=}|=|\ne|\varepsilon} &&\\
&\quad\unsure{\Vert|\text{''}} &&\\[-3pt]
&u_y(g)sx=s\,u_y(g)x\qquad &u_y(g).\,x&=s^{-1}gs\,x\\
&\quad\deleted{\unsure{\Leftarrow|\Uparrow|\Leftrightarrow}\,?} &&=u_x(s^{-1}gs).\,x\\
&\deleted{\bigl(u_x\circ\operatorname{int}(s)\bigr)(g)y=u_x(sgs^{-1})(y)=} &&\\
&u_y(g)=u_x(s^{-1}gs) &&\bigr]
\end{aligned}
\]}
% archive sheet number: 2
